% BOOK-FILE-ID: f-71d7afa99e82
% 付録：本編の式変形を母関数と数列の性質へ広げる

\section{数列をまとめて見る}\label{節：数列をまとめて見る}

本編では,添字をずらして引く操作を何度も使ってきた.
ここではそれを,少し突飛な方法で見直してみよう.
数列の各項を,$x$の累乗の係数として1つの式に並べてしまうのだ.
添字は項の位置を表していたが,その位置を$x$の次数として記録すれば,$x$を掛けるだけで1つ先へずらせる.
等比型$a_0=1$, $a_{n+1}=2a_n$の列$1,2,4,8,\ldots$なら,
\[
  G(x)=1+2x+4x^2+8x^3+\cdots
\]
と書く.ここでは$x$に値を入れて和を求めるのではなく,各次数の係数を記録している.
$2x$を掛けると,一つ先へずれた係数が元の係数と重なる.
\[
  \begin{aligned}
    2xG(x)&=2x+4x^2+8x^3+\cdots,\\
    G(x)-2xG(x)&=1.
  \end{aligned}
\]
したがって$G(x)=\frac{1}{1-2x}$である.「ずらして引く」が,ただの掛け算と引き算になったのだ.
このように各項を係数として並べた式を\bookterm{generating}{母関数}といい,本編の\bookterm{characteristic}{特性方程式}や総和型も同じ仕組みで扱える.

\subsection{母関数と数列の操作}
一般の\bookterm{sequence}{数列}$a_0,a_1,\ldots$では
\[
  A(x)=a_0+a_1x+a_2x^2+a_3x^3+\cdots
  =\sum_{n=0}^{\infty}a_nx^n
\]
とする.収束を仮定せず,各次数の係数を記録する形式的な記号として扱う.項の操作を代数の操作へ移すことが目的である.

たとえば $A(x)$ に $x$ を掛けると,
\[
  xA(x)=a_0x+a_1x^2+a_2x^3+\cdots
\]
となる.
最初に$0$を置き,項を一つずらした形になる.総和型では
$S_n=\sum_{k=0}^n a_k$とし,その\bookterm{generating}{母関数}を$S(x)$と書くと
\[
  S(x)=\frac{A(x)}{1-x}.
\]
これは,形式的な等比級数$\frac{1}{1-x}=1+x+x^2+\cdots$を掛けたとき,$x^n$の係数が$a_0+\cdots+a_n$になるからである.
逆に$(1-x)S(x)=A(x)$では,係数が$S_n-S_{n-1}=a_n$になる.\footnote{ただし$n=0$では$S_{-1}=0$として読む.}
総和型の「ずらして引く」は$1-x$倍になった.$a_1$からの和なら$a_0=0$として使える.

また,二つの\bookterm{sequence}{数列}の項を
\[
  c_n=\sum_{k=0}^{n}a_kb_{n-k}
\]
のように組み合わせる操作は\bookterm{convolution}{畳み込み}と呼ばれ,\bookterm{generating}{母関数}を掛けたときの $x^n$ の係数になる.
例えば全体を二つに分けて数える場合,それぞれの個数を掛けて分け方ごとに足す.このような和を,\bookterm{generating}{母関数}の積で表せる.

\begin{bookpoint}
  \textbf{POINT}\par
  各項を係数として母関数に記録すると,添字のずれは$x$倍,総和は$\frac{1}{1-x}$倍,畳み込みは積になる.
  ここでは級数の収束を仮定せず,各次数の係数を比較する.
\end{bookpoint}

\subsection{分母に現れる漸化式}
この翻訳をフィボナッチ\bookterm{sequence}{数列}で確かめてみよう.
\[
  F_0=0,\quad F_1=1,\qquad F_{n+2}=F_{n+1}+F_n
\]
とし,$F(x)=\sum_{n=0}^{\infty}F_nx^n$ と置く.
左辺の初めの項を取り出すと,
\[
  F(x)-x=\sum_{n=0}^{\infty}F_{n+2}x^{n+2}
\]
である.
\bookterm{recurrence}{漸化式} $F_{n+2}=F_{n+1}+F_n$ を代入すると,右辺の二つの和はそれぞれ $xF(x)$ と $x^2F(x)$ になる.
したがって
\[
  F(x)-x=xF(x)+x^2F(x)
\]
であり,
\[
  F(x)=\frac{x}{1-x-x^2}
\]
を得る.
\bookterm{recurrence}{漸化式}が\bookterm{generating}{母関数}の分母になった.

分母の方程式$1-x-x^2=0$に$r=\frac{1}{x}$を入れると,$r^2-r-1=0$になる.
分母が0になる点は特性方程式の解の逆数だった.行列の\bookterm{eigenvalue}{固有値}が1回の変化で何倍になるかを表すのに対し,母関数の分母は全項の係数関係を表す.

本編の符号に合わせ,$a_{n+2}+pa_{n+1}+qa_n=0$ $(n\ge0)$を考えても同じである.
$A(x)$の係数を比べると
\[
  (1+px+qx^2)A(x)=a_0+(a_1+pa_0)x.
\]
次数$2$以上の係数が\bookterm{recurrence}{漸化式}によってすべて0になるので,初期値の部分だけが右辺に残る.
\bookterm{characteristic}{特性方程式}$\lambda^2+p\lambda+q=0$の根を$\lambda_1,\lambda_2$とすると
\[
  1+px+qx^2=(1-\lambda_1x)(1-\lambda_2x).
\]
異なる0でない根の場合,\bookterm{partial}{部分分数分解}して
\[
  A(x)=\frac{C}{1-\lambda_1x}
  +\frac{D}{1-\lambda_2x}
\]
と書ける.
$\frac{1}{1-\lambda x}=\sum_{n=0}^\infty\lambda^n x^n$の係数を比べれば,$a_n=C\lambda_1^n+D\lambda_2^n$となり,本編の\bookterm{general}{一般項}と同じ形が得られる.
根の一つが0の場合も扱えるが,そのときは多項式の部分を分けて係数を読めばよい.

重解$r\ne0$の場合には分母が$(1-rx)^2$になる.
二つの等比級数を掛けると,$x^n$の係数は
\[
  \sum_{k=0}^n r^k r^{n-k}=(n+1)r^n
\]
だから
\[
  \frac1{(1-rx)^2}
  =\sum_{n=0}^\infty(n+1)r^n x^n.
\]
重解で$n$が現れた理由は,同じ因子の二乗からも読める.

\begin{bookpoint}
  \textbf{POINT}\par
  定数係数の同次線形漸化式は母関数の分母に,初期値は分子に現れる.
  分母の因子に応じて部分分数分解し,係数を読むと一般項に戻せる.重なった因子からは,累乗に掛かる$n$の式が現れる.
\end{bookpoint}

\BookImport{f-73b81d4e9a02}

\subsection{長期的な増え方}
\sectionref*{節：隣接3項間の漸化式を解く}のフィボナッチ数列の式$F_n=\frac{\varphi^n-\psi^n}{\sqrt5}$では,$|\psi|<1<\varphi$なので
\[
  F_n\sim\frac{\varphi^n}{\sqrt5}
\]
となる.
ここで $\sim$ は,二つの量の比が $1$ に近づくという意味である.
大きい絶対値の根が増大を支配する.同じ根を分母から読むことで,\bookterm{generating}{母関数}からも長期的な増大が読み取れる.

\subsection{畳み込みと積}
次は,項どうしを掛けて足す\bookterm{convolution}{畳み込み}を,数え上げの問題で使う.
因子の並びを変えず,二つずつ掛ける括弧の付け方を数えてみよう.
三つの因子$a,b,c$なら,$(ab)c$と$a(bc)$の二通りがある.
$n+1$個の因子に対する個数を$C_n$と書く.
一因子なら掛け算がなく$C_0=1$,二因子なら$C_1=1$,三因子なら$C_2=2$となる.
この列をカタラン数という.各掛け算を左右の分岐で描けば,左右を区別する二分木も数えている.

\paragraphbf{最後の掛け算で分ける}
$n+2$個の因子を掛け終える直前には,左側の積と右側の積が一つずつできている.
左が$k+1$個の因子なら,右は$n-k+1$個である.
左右それぞれの括弧の付け方を選ぶと,$C_kC_{n-k}$通りになる.
$k$は$0$から$n$まで動き,この分け方は最後の掛け算から一意に決まる.
従って,
\[
  C_0=1,\qquad C_{n+1}=\sum_{k=0}^{n}C_kC_{n-k}
\]
となる.左右を選ぶ積を全分け方で足すので,畳み込みになる.
\bookterm{generating}{母関数}$C(x)=\sum_{n=0}^{\infty}C_nx^n$では
\[
  C(x)=1+xC(x)^2
\]
という二次方程式になる.これを解くと
\[
  C(x)=\frac{1-\sqrt{1-4x}}{2x}
\]
である.
ここで平方根は,定数項が$1$となる形式的な級数として読む.\footnote{$H(x)=1+h_1x+h_2x^2+\cdots$とし,$H(x)^2=1-4x$の係数を合わせる.$x^n$の係数には$2h_n$と,先に決まった係数だけが現れるので,$h_n$を順に一意に求められる.ここでは数値を代入した級数の収束を使っていない.}
逆の符号では$x^{-1}$の項が生じるため,$C_0=1$から始まる母関数には合わない.
$x$の次数ごとに係数を求めると,
\[
  C(x)=1+x+2x^2+5x^3+\cdots
\]
となり,三因子の二通り,四因子の五通りも係数に戻る.

\paragraphbf{係数を取り出す}
では,$C_n$を$n$の式で表せないだろうか.
平方根を展開するのは高校範囲を超えるので,ここでは別の道を通ろう.
この計算では形式的な級数の微分を使うので,初めて読むときは飛ばしてよい.
脚注の$H(x)=\sqrt{1-4x}=\sum_{n=0}^{\infty}h_nx^n$を使う.$h_0=1,\ h_1=-2$である.
$H(x)^2=1-4x$の両辺を項ごとに微分すると$2H(x)H'(x)=-4$である.
両辺に$H(x)$を掛けて$H(x)^2=1-4x$を使えば,
\[
  (1-4x)H'(x)=-2H(x)
\]
が成り立つ.\footnote{級数を多項式と同じように項ごとに微分している.係数の並びとして扱う形式的な級数の立場では,収束を考えずにこの計算を正当化できる.}
両辺の$x^n$の係数を比べると,
\begin{align*}
  &(n+1)h_{n+1}-4nh_n=-2h_n\\
  &\therefore\quad h_{n+1}=\frac{4n-2}{n+1}h_n
\end{align*}
となる.なんと\sectionref*{節：階比型}である.
$C(x)=\frac{1-H(x)}{2x}$より$C_n=-\frac{h_{n+1}}{2}$なので,$C_n$も
\[
  C_n=\frac{2(2n-1)}{n+1}C_{n-1}\qquad(n\ge1)
\]
という階比型を満たす.書き下していけば
\begin{align*}
  C_n&=\prod_{k=1}^{n}\frac{2(2k-1)}{k+1}\\
  &=\frac{2^n\cdot1\cdot3\cdots(2n-1)}{(n+1)!}\\
  &=\frac{(2n)!}{n!\,(n+1)!}
\end{align*}
すなわち
\[
  C_n=\frac{1}{n+1}\binom{2n}{n}
\]
を得る.\footnote{最後の変形では$(2n)!=\{1\cdot3\cdots(2n-1)\}\cdot\{2\cdot4\cdots(2n)\}=1\cdot3\cdots(2n-1)\cdot2^nn!$を使った.}
総和を含む非線形の漸化式が,母関数を経由すると,本編の最初に学んだ階比型に帰着したのだ.
母関数は,扱える漸化式の範囲を大きく広げてくれる道具なのである.

\bookterm{nonlinear}{非線形}でも,和と積が母関数の操作に合えば代数方程式へ移せる.すべての非線形の漸化式に使えるわけではない.

この発想は,\bookterm{chebyshev}{チェビシェフ多項式}にも戻ってくる.
その\bookterm{recurrence}{漸化式}に\bookterm{generating}{母関数}を作ると,
\[
  \sum_{n=0}^{\infty}T_n(x)t^n
  =\frac{1-xt}{1-2xt+t^2}
\]
となる.
実際,$G(t)=\sum_{n=0}^\infty T_n(x)t^n$とし,$x$を定数として扱えば
\[
  (1-2xt+t^2)G(t)=1-xt.
\]
右辺は$T_0(x)=1$, $T_1(x)=x$に由来し,二次以上の係数は\bookterm{recurrence}{漸化式}によって打ち消される.
余弦の整数倍を表す列が,一つの有理関数にもまとまった.\bookterm{recurrence}{漸化式},\bookterm{trigonometric}{三角関数},\bookterm{generating}{母関数}は同じ列の三つの表現である.

\paragraphbf{別の広がり：階乗の番号を広げる}
この部分では広義積分など数学IIIの内容を使う.初めて読むときは飛ばしてもよい.
母関数は各項を係数に記録した.次は,項の\textbf{番号そのもの}を整数以外へ広げる.
階乗は$u_0=1$, $u_{n+1}=(n+1)u_n$という\bookterm{ratio}{階比}型の例である（\booknavref{解法：階比型：f(n)が1次式の積の場合}{1次式の積の場合}）.
その関係を広げるガンマ関数は,正の$x$に対して
\[
  \Gamma(x)=\int_0^\infty t^{x-1}e^{-t}\,dt
\]
と定められる.
部分積分をすると
\[
  \Gamma(x+1)=x\Gamma(x)
\]
が現れる.最初の値も
\[
  \Gamma(1)=\int_0^\infty e^{-t}\,dt=1
\]
なので,整数$n\ge0$では$\Gamma(n+1)=n!$となる.
番号を広げた階乗が積分や確率につながる.\bookterm{recurrence}{漸化式}から関数を作る別の例である.

\subsection{確率と整数への応用}
確率に目を向けると, 状態の移り変わりも\bookterm{recurrence}{漸化式}で記録できる.
晴れ・雨の確率を並べた\bookterm{column-vector}{列ベクトル}を$\boldsymbol p_n$とする.晴れ,雨の翌日に晴れる確率がそれぞれ$0.8,0.3$なら
\[
  \boldsymbol p_{n+1}
  =
  \begin{pmatrix}0.8&0.3\\0.2&0.7\end{pmatrix}
  \boldsymbol p_n
  =P\boldsymbol p_n
\]
と書ける.
各列の和は$1$なので,確率の合計は$1$に保たれる.
晴れの確率を$u_n$とすれば,雨の確率は$1-u_n$だから
\[
  u_{n+1}=0.8u_n+0.3(1-u_n)=0.5u_n+0.3.
\]
\sectionref*{節：特性方程式型}で扱った定数項をもつ式そのもので,\bookterm{fixed}{不動点}は$0.6$である.
したがって
\[
  u_n-0.6=(u_0-0.6)(0.5)^n
\]
となり,晴れ・雨の確率は$(0.6,0.4)$へ近づく.
\bookterm{matrix}{行列}では,$(1,1)P=(1,1)$が確率の合計を保つことを表し,\bookterm{eigenvalue}{固有値}$0.5$は\bookterm{fixed}{不動点}からのずれが1回ごとに何倍になるかを表す.
不動点と等比型が,確率の長期予測にも使えた.

数論では, ユークリッドの互除法が余りを次々に作る漸化式として表せる.
二つの正の整数$a_0,a_1$から始める.
$a_{n-1}$ を $a_n$ で割った商を $q_n$, 余りを $a_{n+1}$ とすれば,
\[
  a_{n-1}=q_na_n+a_{n+1},
  \qquad 0\le a_{n+1}<a_n
\]
となる.
余りが0なら計算を止め,そうでなければ2つの数を$(a_n,a_{n+1})$へ置き換える.
$a_{n+1}=a_{n-1}-q_na_n$と逆向きの式から,共通の約数は変わらない.従って
\[
  \gcd(a_{n-1},a_n)=\gcd(a_n,a_{n+1}).
\]
正の余りは小さくなるので有限回で止まる.
余り$a_{n+1}$が$0$になったときの除数$a_n$,すなわち最後の0でない項が最大公約数となる.
確率の合計と同じく,変わらない量から計算の意味を説明できる.

整数係数・整数初期値の隣接3項型では,各項を整数$m\ge2$で割った余りだけを残してみよう.
状態 $(a_n,a_{n+1})$ の各成分は $0,1,\ldots,m-1$ のいずれかなので, 状態は高々 $m^2$ 通りしかない.
次の状態が今の状態だけで決まるなら, いつか同じ状態が現れ, その後は同じ並びを繰り返す.
したがって,余りの列はある項以降で周期的になる.
たとえばフィボナッチ\bookterm{sequence}{数列}を $2$ で割った余りは
\[
  0,1,1,0,1,1,0,1,1,\ldots
\]
と周期 $3$ で最初から繰り返す.
一方,漸化式を逆向きにたどれない場合には,最初から周期的になるとは限らない.
たとえば $a_1=a_2=1$,$a_{n+2}=2a_{n+1}$ を $4$ で割った余りは
\[
  1,1,2,0,0,0,\ldots
\]
となり,$a_4$ 以降は周期 $1$ で繰り返すが,初めの $1,1,2$ は繰り返されない.
同じ漸化式でも,実数の増大と余りの周期では見える性質が変わる.

本編の\bookterm{trigonometric}{三角関数}型に似た二次式でも,同じ解法が使えるとは限らない.それでも増大は評価できる.
たとえば
\[
  a_0=0,\qquad a_{n+1}=a_n^2+1
\]
では $0,1,2,5,26,\ldots$ と増える.
$n\ge2$ では $a_n\ge2$ であり, $a_{n+1}\ge a_n^2$ だから, 帰納的に
\[
  a_n\ge 2^{\,2^{n-2}}
\]
となる.
一般項を求めずに,急速な増大を評価できた.
また, \bookterm{sequence}{数列}が単調で有界と分かれば収束し, 漸化式が $a_{n+1}=f(a_n)$ で $f$ が連続なら, \bookterm{limit}{極限} $L$ は $L=f(L)$ を満たす.
単調性, 有界性, \bookterm{fixed}{不動点}, 増加速度は, \bookterm{general}{一般項}とは別の問いに答えてくれる.

行列・差分・母関数は,本編の操作を別の表現へ移す道具だった.
不動点からのずれ,\bookterm{left-eigen}{左固有ベクトル}で取り出す量,添字のずれは,それぞれ同じ漸化式の別の表し方である.
一般項に加え,増大・確率・周期・保存される量も,見たい性質に合わせて調べられる.

\begin{bookpoint}
  \textbf{POINT}\par
  調べたい性質に合わせ,数列全体を母関数で表したり,確率や余りなどの量に注目したりする.
  極限の候補を不動点の式から得た場合も,収束そのものは単調性・有界性などで別に確かめる.
\end{bookpoint}
